\begin{tikzpicture}

    % ============================================================
    % ÉCHELLES ET PARAMÈTRES
    % ============================================================

    \pgfmathsetmacro{\echelleY}{1.5}
    \pgfmathsetmacro{\echelleX}{1.25}

    \def\lam{4}           % Longueur d'onde
    \def\amp{1}           % Amplitude

    % Deux instants :
    % Le profil se déplace vers les x croissants
    \def\xUn{1.5}
    \def\DeltaX{2.5}
    \pgfmathsetmacro{\xDeux}{\xUn+\DeltaX}

    % Phase choisie pour avoir un point P identifiable
    \def\phi{-90}

    % ============================================================
    % CLIP
    % ============================================================

    \clip
        (-1.5*\echelleX,-2.3*\echelleY)
        rectangle
        (11.0*\echelleX,2.3*\echelleY);

    % ============================================================
    % GRILLE
    % ============================================================

    \backgrowngrid
        {-1}{-2}{10}{2}
        {\echelleX}{\echelleY}
        {\colorThree!50!\colorTwo}
        {\colorTwo}
        {0.5ex}

    % ============================================================
    % ÉQUATION
    % ============================================================

    \node[
        fill=\colorThree!25!\colorTwo,
        rounded corners=6pt,
        inner sep=8pt,
        font=\sffamily\bfseries\fontsize{18pt}{22pt}\selectfont,
        text=\colorOne
    ]
    at (4.5*\echelleX,1.85*\echelleY)
    {
        $\bm{y(x,t)=f(x-ct)}$
    };

    % ============================================================
    % PROFIL À t_1
    % ============================================================

    \draw[
        line width=0.8ex,
        color=\colorFour,
        domain=-1:10,
        samples=250,
        smooth
    ]
    plot
    (
        \x*\echelleX,
        {
            \echelleY*\amp*
            cos(360/\lam*\x+\phi)
        }
    );

    % ============================================================
    % PROFIL À t_2
    %
    % f(x-c t_2) = f((x-c t_1)-Delta x)
    %
    % Donc translation vers la DROITE
    % ============================================================

    \draw[
        line width=0.8ex,
        color=\colorSix,
        domain=-1:10,
        samples=250,
        smooth
    ]
    plot
    (
        \x*\echelleX,
        {
            \echelleY*\amp*
            cos(360/\lam*(\x-\DeltaX)+\phi)
        }
    );

    % ============================================================
    % POSITION DES POINTS IDENTIFIABLES
    % ============================================================

    \pgfmathsetmacro{\xPUn}{\xUn*\echelleX}
    \pgfmathsetmacro{\xPDeux}{\xDeux*\echelleX}

    \pgfmathsetmacro{\yP}{
        \echelleY*\amp*
        cos(360/\lam*\xUn+\phi)
    }

    % Point P à t_1
    \fill[
        color=\colorFour
    ]
    (\xPUn,\yP)
    circle (0.14cm);

    \node[
        above left,
        font=\sffamily\bfseries\fontsize{15pt}{18pt}\selectfont,
        text=\colorFour
    ]
    at (\xPUn,\yP)
    {$P_1$};

    % Point P à t_2
    \fill[
        color=\colorSix
    ]
    (\xPDeux,\yP)
    circle (0.14cm);

    \node[
        above right,
        font=\sffamily\bfseries\fontsize{15pt}{18pt}\selectfont,
        text=\colorSix
    ]
    at (\xPDeux,\yP)
    {$P_2$};

    % ============================================================
    % PROJECTIONS DES POINTS SUR L'AXE
    % ============================================================

    \draw[
        dashed,
        line width=0.45ex,
        color=\colorFour
    ]
    (\xPUn,\yP)
    -- (\xPUn,0);

    \draw[
        dashed,
        line width=0.45ex,
        color=\colorSix
    ]
    (\xPDeux,\yP)
    -- (\xPDeux,0);

    % ============================================================
    % DÉPLACEMENT DU POINT
    % ============================================================

    \draw[
        ->,
        >=Latex,
        line width=0.8ex,
        color=\colorOne
    ]
    (\xPUn,0.85*\echelleY)
    --
    node[
        midway,
        above,
        fill=\colorTwo,
        inner sep=3pt,
        font=\sffamily\bfseries\fontsize{14pt}{17pt}\selectfont,
        text=\colorOne
    ]
    {
        $\Delta x=c(t_2-t_1)$
    }
    (\xPDeux,0.85*\echelleY);

    % ============================================================
    % LÉGENDE DES DEUX PROFILS
    % ============================================================

    \draw[
        line width=0.8ex,
        color=\colorFour
    ]
    (0.2*\echelleX,-1.55*\echelleY)
    --
    (1.0*\echelleX,-1.55*\echelleY)
    node[
        right,
        text=\colorFour,
        font=\sffamily\bfseries\fontsize{13pt}{16pt}\selectfont
    ]
    {$t=t_1$};

    \draw[
        line width=0.8ex,
        color=\colorSix
    ]
    (3.2*\echelleX,-1.55*\echelleY)
    --
    (4.0*\echelleX,-1.55*\echelleY)
    node[
        right,
        text=\colorSix,
        font=\sffamily\bfseries\fontsize{13pt}{16pt}\selectfont
    ]
    {$t=t_2>t_1$};

    % ============================================================
    % AXE X
    % ============================================================

    \draw
    (-1.3*\echelleX,0)
    edge[
        thick,
        line width=0.8ex,
        ->,
        >=Triangle,
        color=\colorOne
    ]
    node[
        pos=1,
        right,
        font=\sffamily\bfseries\fontsize{16pt}{18pt}\selectfont,
        text=\colorOne
    ]
    {$x$}
    (10.8*\echelleX,0);

    % ============================================================
    % AXE Y
    % ============================================================

    \draw
    (0,-2.0*\echelleY)
    edge[
        thick,
        line width=0.8ex,
        ->,
        >=Triangle,
        color=\colorOne
    ]
    node[
        pos=1,
        above,
        font=\sffamily\bfseries\fontsize{16pt}{18pt}\selectfont,
        text=\colorOne
    ]
    {$y$}
    (0,2.1*\echelleY);

    % ============================================================
    % CONCLUSION GRAPHIQUE
    % ============================================================

    \node[
        fill=\colorFour!15!\colorTwo,
        rounded corners=5pt,
        inner sep=7pt,
        font=\sffamily\bfseries\fontsize{15pt}{18pt}\selectfont,
        text=\colorFour
    ]
    at (7.0*\echelleX,-1.55*\echelleY)
    {
        propagation vers les $\bm{x}$ croissants
    };

\end{tikzpicture}